\documentclass[11pt]{article}
\usepackage[margin=1in]{geometry}
\usepackage{amsmath,amssymb,amsthm}
\usepackage[hidelinks]{hyperref}
\newtheorem{proposition}{Proposition}
\newtheorem{lemma}[proposition]{Lemma}
\newtheorem{corollary}[proposition]{Corollary}
\theoremstyle{remark}
\newtheorem{remark}[proposition]{Remark}
\newcommand{\C}{\mathbb C}
\newcommand{\HP}{\mathcal H_\Psi}
\newcommand{\HF}{\mathcal H_\Phi}
\newcommand{\K}{\mathcal K}
\newcommand{\Ran}{\operatorname{Ran}}
\newcommand{\kr}[2]{[#1,#2]}
\newcommand{\ket}[1]{\lvert #1\rangle}
\newcommand{\bra}[1]{\langle #1\rvert}
\title{The shift example as a BRST model}
\author{Working note}
\date{}
\begin{document}
\maketitle

\begin{abstract}
\noindent
We test a BRST-type reading of the polar correspondence. From the polar part of a preparation we build a nilpotent operator that is self-adjoint for the Krein form. Its cohomology is the kernel of the preparation together with its defect: it records what the specified presentation leaves unpaired. A sector-preserving evolution acts on this quotient exactly when it commutes with the charge, and that holds exactly when the polar map intertwines the dynamics. Requiring the observer's evolution to act on the quotient is a physical premise. Choosing the Krein form does not impose it. In the shift example of the consolidated manuscript (Section~7) the reduced space is the missing line. Under the premise, for every excitation energy other than one, the missing line lies in the common defect of every family whose joint charge the evolution preserves. The manuscript imposes this spectral exclusion as a restriction; here it follows from the premise. The premise constrains free evolution only. Excluding transfers into the missing line by the observer's interventions needs a second premise: every admitted intervention also acts on the quotient. Every matched measurement acts trivially on the reduced space, and the reduced space is orthogonal to the luminous ground. Under R4 it is zero when there is no excess. The reduction removes each matched event from the count, whereas identifying the two descriptions of an event would keep one copy. We read the reduced space as a record of what remains unpaired by the specified presentation, not as an inventory of real appearances.
\end{abstract}

\section{The proposal and its one departure}

The consolidated manuscript computes probabilities with the positive Hilbert inner product in each sector. Its indefinite form only labels the sectors (Section~3). The proposal tested here uses the indefinite form as the inner product of a reduced theory, as in BRST quantization: states are reduced by a nilpotent charge, and the reduced space is its cohomology with the induced form. That is the one assumption that departs from the manuscript. Two premises are physical. The observer's evolution must act on the quotient (Lemma~\ref{lem:quotient}). Where interventions are discussed, every admitted intervention must act on it as well (Section~\ref{sec:shift}). How the reduced space is read is a separate, interpretive matter (Section~\ref{sec:reading}). Everything else below is a statement about operators, and each finite-dimensional claim is checked in \texttt{test\_brst\_shift.py}.

\section{The charge of a preparation}

Let $\K=\HP\oplus\HF$ with Hilbert product $\langle z,w\rangle$, linear in the second argument, fundamental symmetry $J(a\oplus b)=a\oplus(-b)$, Krein form $\kr zw=\langle z,Jw\rangle$ and Krein adjoint $X^\sharp=JX^*J$. Let $A:\HP\to\HF$ be a bounded preparation with polar decomposition $A=V|A|$. The partial isometry $V$ has initial space $S=(\ker A)^\perp$ and final space $D=\overline{\Ran A}$. Only $V$ enters below, so a closed unbounded $A$ can be treated the same way. The \emph{aligned graph} and its orthogonal projection are
\begin{equation}
 G=\{\psi\oplus V\psi:\psi\in S\},\qquad
 P_G=\tfrac12\begin{pmatrix}P_S&V^*\\ V&P_D\end{pmatrix},
 \label{eq:graph-projection}
\end{equation}
and the \emph{charge} of $A$ is
\begin{equation}
 Q_A=P_GJ=\tfrac12\begin{pmatrix}P_S&-V^*\\ V&-P_D\end{pmatrix}.
 \label{eq:charge}
\end{equation}

\begin{proposition}\label{prop:charge}
\begin{enumerate}
\item[(a)] $Q_A^2=0$ and $Q_A^\sharp=Q_A$.
\item[(b)] $\Ran Q_A=G$.
\item[(c)] $\ker Q_A=G\oplus(\ker A\oplus D^\perp)$, an orthogonal sum.
\item[(d)] The Krein form vanishes between $\Ran Q_A$ and $\ker Q_A$, so it induces a form on $H(Q_A)=\ker Q_A/\Ran Q_A$. The map $\ker A\oplus D^\perp\to H(Q_A)$ is a linear bijection, and the induced form is $\kr{k\oplus d}{k'\oplus d'}=\langle k,k'\rangle-\langle d,d'\rangle$. It is positive on $\ker A$ and negative on $D^\perp$.
\end{enumerate}
\end{proposition}

\begin{proof}
(a) $J$ is unitary and self-adjoint and carries $G$ onto $G_-=\{\psi\oplus(-V\psi)\}$, so $JP_GJ=P_{G_-}$. The two graphs are orthogonal, since $\langle\psi\oplus V\psi,\varphi\oplus(-V\varphi)\rangle=\langle\psi,\varphi\rangle-\langle V\psi,V\varphi\rangle=0$. Hence $Q_A^2=P_G(JP_GJ)=P_GP_{G_-}=0$, and $Q_A^\sharp=J(P_GJ)^*J=JJP_GJ=Q_A$. (b) $J$ is bijective. (c) $Q_Az=0$ exactly when $Jz\perp G$. For $z=a\oplus b$ and $\psi\in S$, $\langle\psi\oplus V\psi,a\oplus(-b)\rangle=\langle\psi,a-V^*b\rangle$, so the condition is $P_Sa=V^*b$, equivalently $P_Db=VP_Sa$. Then $z=(P_Sa\oplus VP_Sa)+((I-P_S)a\oplus(I-P_D)b)$, and the two summands are orthogonal. Conversely $G=\Ran Q_A\subseteq\ker Q_A$, and $Q_A(k\oplus d)=\tfrac12\bigl((P_Sk-V^*d)\oplus(Vk-P_Dd)\bigr)=0$ for $k\in\ker A$, $d\in D^\perp$. (d) $\kr{Q_Az}w=\kr z{Q_Aw}=0$ for $w\in\ker Q_A$. By (b) and (c), $\ker A\oplus D^\perp$ is a complement of $\Ran Q_A$ in $\ker Q_A$, and representatives differ by elements on which the form vanishes.
\end{proof}

\begin{remark}[Witt pairs]
For a unit vector $\psi\in S$ put $n_\psi=(\psi\oplus V\psi)/\sqrt2$ and $m_\psi=Jn_\psi=(\psi\oplus(-V\psi))/\sqrt2$. Both are Krein-null, $\kr{n_\psi}{m_\psi}=1$, $Q_An_\psi=0$ and $Q_Am_\psi=n_\psi$. The aligned vector is exact, and its anti-aligned companion is the non-closed partner. Such pairs span $S\oplus D=G\oplus G_-$, so every matched physical state cancels together with its presentation: the matched part of $\K$ contributes nothing to $H(Q_A)$. These are the graph vectors $\mathcal N_0$ and $\mathcal N_\pi$ of the manuscript (Section~3). Identifying the two descriptions $s\oplus0$ and $0\oplus Vs$ of one event is a different operation. It quotients $S\oplus D$ by $G_-$ and leaves one copy of each event, isomorphic to $G\cong S$, on which the paired effect $E\oplus VEV^*$ has the matched probability $\langle s,Es\rangle$. Cohomology leaves zero copies.
\end{remark}

\section{Gauge-invariant operators}

Call a bounded operator \emph{invariant} if it commutes with $Q_A$. Invariant operators preserve $\ker Q_A$ and $\Ran Q_A$ and therefore act on $H(Q_A)$. We treat even operators $X=X_\Psi\oplus X_\Phi$, which commute with $J$. These are the operators compatible with the superselection rule of the manuscript's Section~3.

\begin{lemma}\label{lem:invariant}
$X=X_\Psi\oplus X_\Phi$ is invariant if and only if
\[
 [X_\Psi,P_S]=0,\qquad [X_\Phi,P_D]=0,\qquad X_\Phi V=VX_\Psi,\qquad X_\Psi V^*=V^*X_\Phi .
\]
Its action on $H(Q_A)\cong\ker A\oplus D^\perp$ is then $X_\Psi|_{\ker A}\oplus X_\Phi|_{D^\perp}$.
\end{lemma}

\begin{proof}
Since $[X,J]=0$, $[X,Q_A]=[X,P_G]J$, which vanishes exactly when $[X,P_G]=0$. Comparing the four blocks of $XP_G$ and $P_GX$ from Eq.~\eqref{eq:graph-projection} gives the four conditions. They imply that $X$ preserves $\ker A\oplus D^\perp$, which by Proposition~\ref{prop:charge} represents $H(Q_A)$.
\end{proof}

For self-adjoint $X$ the fourth condition is the adjoint of the third. On $S\oplus D$ an invariant even operator is therefore determined by $X_\Psi|_S$, through $X_\Phi|_D=VX_\Psi|_SV^*$.

\begin{lemma}[Evolution on the quotient]\label{lem:quotient}
A unitary group $U(t)=U_\Psi(t)\oplus U_\Phi(t)$ acts on $H(Q_A)$, in the sense that $U(t)\Ran Q_A\subseteq\Ran Q_A$ for every real $t$, exactly when $U(t)$ commutes with $Q_A$ for every $t$.
\end{lemma}

\begin{proof}
If $U(t)G\subseteq G$ for all $t$, applying this at $-t$ gives $U(t)G=G$. A unitary that maps a closed subspace onto itself also maps its orthogonal complement onto itself, so $U(t)$ commutes with $P_G$. It commutes with $J$, hence with $Q_A=P_GJ$. Conversely, an operator commuting with $Q_A$ preserves $\Ran Q_A$ and $\ker Q_A$.
\end{proof}

Requiring the observer's evolution to act on the quotient is therefore the same as requiring it to commute with $Q_A$. This is a physical premise of the reading. The choice of the Krein form as an inner product does not by itself impose it. The premise concerns free evolution. Interventions need a premise of their own (Section~\ref{sec:shift}).

\begin{corollary}\label{cor:consequences}
\begin{enumerate}
\item[(a)] \emph{Matched effects.} For an effect $E=P_SEP_S$ on $S$, the paired effect $E\oplus VEV^*$ is invariant and acts as zero on $H(Q_A)$.
\item[(b)] \emph{Dynamics.} Let $H_\Psi$ and $H_\Phi$ be self-adjoint, possibly unbounded. The unitary group of $H=H_\Psi\oplus H_\Phi$ commutes with $Q_A$ for every $t$ exactly when $S$ reduces $H_\Psi$, $D$ reduces $H_\Phi$, and the polar map intertwines the groups, $e^{-itH_\Phi}V=Ve^{-itH_\Psi}$ for every $t$. The evolution then acts on $H(Q_A)$ as $e^{-itH_\Psi}|_{\ker A}\oplus e^{-itH_\Phi}|_{D^\perp}$.
\item[(c)] \emph{No transfer.} An invariant operator maps $\Ran Q_A$ into itself. In particular it never sends an exact vector to one with a nonzero class in $H(Q_A)$.
\end{enumerate}
\end{corollary}

\begin{proof}
(a) Lemma~\ref{lem:invariant} with $X_\Psi=E$ and $X_\Phi=VEV^*$: the action on $H(Q_A)$ is $E|_{\ker A}\oplus VEV^*|_{D^\perp}=0$. (b) Lemma~\ref{lem:invariant} applied to each $U(t)$. The fourth condition at $t$ is the adjoint of the third at $-t$, and commuting with $P_S$ and $P_D$ for every $t$ means that $S$ and $D$ reduce the generators. (c) $X Q_Az=Q_AXz$.
\end{proof}

It is the polar map that must intertwine, not $A$ itself. A positive invertible $A$ on a single space has polar map $V=I$, which intertwines any dynamics with itself, while $|A|=A$ can mix energies.

\begin{proposition}[Spectral obstruction]\label{prop:obstruction}
Let $B:\mathcal H'_\Psi\to\HF$ be bounded with polar part $W$ and final space $D_B$, and let the unitary group of $H'=H'_\Psi\oplus H_\Phi$ commute with $Q_B$. If $\lambda$ is an eigenvalue of $H_\Phi$ but not of $H'_\Psi$, every eigenvector of $H_\Phi$ for $\lambda$ lies in the defect $D_B^\perp$.
\end{proposition}

\begin{proof}
By Lemma~\ref{lem:invariant} applied to each group element, $e^{-itH'_\Psi}W^*=W^*e^{-itH_\Phi}$ for every $t$. For $H_\Phi h=\lambda h$ this gives $e^{-itH'_\Psi}W^*h=e^{-i\lambda t}W^*h$ for every $t$. So $W^*h$ is an eigenvector of $H'_\Psi$ for $\lambda$ unless it is zero. Hence $W^*h=0$, that is, $h\perp D_B$.
\end{proof}

A family of preparations $A_c:\HP\to\HF$ enters through its joint map $B=(A_c)_c$ on $\bigoplus_c\HP$. The defect of $B$ is the common defect $D_{\mathcal C}^\perp$ of the manuscript's Eq.~(59). By Proposition~\ref{prop:charge} applied to $B$, $H(Q_B)\cong\ker B\oplus D_{\mathcal C}^\perp$. The kernel collects redundant presentations, combinations of inputs with the same total presentation. It enters with positive sign, and the excess enters with negative sign. So whenever presentations overlap and an excess remains, the reduced space is indefinite. Treating $\ker B$ as gauge too, by restricting $B$ to $(\ker B)^\perp$, leaves exactly $D_{\mathcal C}^\perp$.

\paragraph{A ghost grading.} The sector sign $J$ does not grade $Q_A$, which is neither even nor odd under $J$. A grading exists all the same. The manuscript's polar reflection $\mathcal R$ (Section~3), extended by zero off the active carrier, equals $P_G-P_{G_-}$. Then $N=\tfrac12\mathcal R$ satisfies
\[
 [N,Q_A]=Q_A,\qquad N^\sharp=-N ,
\]
with value $+\tfrac12$ on exact vectors, $-\tfrac12$ on their non-closed partners and $0$ on the unpaired directions. These are the defining properties of a ghost-number operator. Every class that survives has ghost number zero, since $N$ vanishes on the unpaired directions. On a graph state of relative phase $\varphi$ the mean of $\mathcal R$ is $\cos\varphi$. The observable that detects the relative phase is therefore twice the ghost number, and it does not commute with $Q_A$. The relative phase is gauge only if $\mathcal R$ is excluded from the observables of the reduced theory. The manuscript admits it as a structural operator, and the BRST reading excludes it.

\section{The shift}
\label{sec:shift}

Take the manuscript's shift example (Section~7.3): orthonormal bases $f_n$ of $\HP$ and $g_n$ of $\HF$, $n\ge0$, and $Af_n=g_{n+1}$, with
\[
 H_\Psi=I-P_{f_0},\qquad H_\Phi^{(\varepsilon)}=I-P_{g_1}+(\varepsilon-1)P_{g_0},\qquad \varepsilon>0 .
\]
The input ground is $x=f_0$ and the output ground is $y=g_1=Ax$, so R3 holds with $a=1$. $A$ is an isometry, so $V=A$, $S=\HP$ (R4) and $D^\perp=\C g_0$ (R5).

\paragraph{Physical space and dynamics.} By Proposition~\ref{prop:charge}, $H(Q_A)=\C g_0$ with Krein norm $-1$. Its sign is a convention, and with the opposite convention the induced norm is the ordinary one. For every $\varepsilon>0$, $H_\Phi^{(\varepsilon)}A=AH_\Psi$ and $g_0$ is an eigenvector. By Corollary~\ref{cor:consequences}(b), the dynamics is invariant, and the reduced space evolves alone, by $e^{-i\varepsilon t}$. The exact vectors are spanned by the aligned pairs $(f_n\oplus g_{n+1})/\sqrt2$.

\paragraph{Which presentations are consistent.} The rotated shifts $A^{(\alpha)}$ of the manuscript's Eq.~(92) are isometries with $\|H_\Phi^{(\varepsilon)}A^{(\alpha)}-A^{(\alpha)}H_\Psi\|=|\varepsilon-1||\sin\alpha|$ (Eq.~(93)). By Corollary~\ref{cor:consequences}(b), $Q_{A^{(\alpha)}}$ commutes with the dynamics exactly when $\varepsilon=1$ or $\sin\alpha=0$.

For families, $H'_\Psi=\bigoplus_cH_\Psi$ has eigenvalues $0$ and $1$ only. For $\varepsilon\ne1$, $\varepsilon$ is an eigenvalue of $H_\Phi^{(\varepsilon)}$ but not of $H'_\Psi$, with eigenvector $g_0$. Proposition~\ref{prop:obstruction}, applied to the joint map or to each member (Remark~\ref{rem:members} below), gives:

\begin{corollary}\label{cor:shift}
Assume the evolution acts on the quotient. For $\varepsilon\ne1$, every family of presentations whose joint charge commutes with the dynamics has $g_0$ in its common defect. So does every family each of whose members' charges commutes with the dynamics. After redundant presentations are treated as gauge, the reduced space of such a family contains the missing line. At $\varepsilon=1$ the family $\{A,A^{(\alpha)}\}$ with $\sin\alpha\ne0$ is consistent and has no common defect. Both maps send $f_0$ to $g_1$, so $(f_0,-f_0)$ is a redundant presentation, and the reduced space of the single joint charge is not zero: it consists of such redundancies. Only after redundant presentations are treated as gauge is the reduced space zero.
\end{corollary}

\begin{remark}[Members and families]\label{rem:members}
For a family, the premise can be imposed on the quotient of each member or on the quotient of the joint map. The two conditions are independent. Let $H=\operatorname{diag}(0,1,2)$ act on inputs and outputs, $A_1=I$ and $A_2=1\oplus\left(\begin{smallmatrix}2&1\\1&2\end{smallmatrix}\right)$. The two preparations are injective, with the same ground, the same range and the same polar map $I$, so each member's charge commutes with the dynamics. But $\ker B=\{(-A_2\psi)\oplus\psi\}$ is not invariant under $H\oplus H$, because $A_2$ mixes the levels $1$ and $2$, so $Q_B$ does not commute with the dynamics. Conversely, let $H=\operatorname{diag}(0,1)$ act on two input copies of $\C^2$ and on the output $\C^2$. Take $u_0=(0.9,0,\sqrt{0.19},0)$ and $u_1=(0,0.3,0,\sqrt{0.91})$, of input energies $0$ and $1$, and $B=W|B|$ with $Wu_j=g_j$, $|B|u_0=u_0+\tfrac12u_1$, $|B|u_1=\tfrac12u_0+u_1$ and $|B|=0$ on the complement. The joint charge commutes with the dynamics, but neither member's polar map is diagonal, so neither member's charge does. Proposition~\ref{prop:obstruction} applies under either condition. Applied to each member, it places every eigenvector of $H_\Phi$ for an eigenvalue absent from the input spectrum in each member's defect, and hence in the common defect. Only under the joint condition does the evolution act on the family's reduced space $H(Q_B)$ (Lemma~\ref{lem:quotient}).
\end{remark}

The manuscript obtains $\C g_0$ as the common defect by \emph{assuming} that admissible presentations intertwine the dynamics (Section~7.3). Here intertwining by the polar map follows from the premise that the evolution acts on the quotient, and the exclusion of $g_0$ follows in turn. A presentation that breaks the exclusion, such as the map $f_1\mapsto g_0$ of the covariant measure-and-prepare channel of Sections~7.4--7.5, makes the joint charge non-invariant.

\paragraph{Energy transfer.} The operator $K=0\oplus\ket{g_0}\!\bra{g_1}$ satisfies $[H_\Phi^{(\varepsilon)},K_\Phi]=\varepsilon K_\Phi$, so the single-outcome operation $\sigma\mapsto K\sigma K^*$ is time-translation covariant (Section~7.4). It is not invariant. It sends the exact ground event $z_+=(f_0\oplus g_1)/\sqrt2$ to $(0\oplus g_0)/\sqrt2$, whose class in $H(Q_A)$ is nonzero. By Corollary~\ref{cor:consequences}(c), no invariant operator can do that. Time covariance can therefore permit a transfer that charge invariance forbids. The premise on evolution does not exclude this operation: the shift dynamics commutes with $Q_A$ for every $\varepsilon$, and $K$ is available all the same. To conclude that the observer's interventions transfer no weight from exact vectors into the missing line, the model needs a second premise: every admitted operation acts on the quotient, in the sense that each of its Kraus operators maps $\Ran Q_A$ and $\ker Q_A$ into themselves. The condition does not depend on the Kraus representation, because the Kraus operators of one representation are linear combinations of those of any other. Invariant operators satisfy it (Corollary~\ref{cor:consequences}(c)). Under the second premise an exact vector stays exact, and no admitted operation sends weight from an exact event into a nonzero class. The first premise does not imply the second, and the choice of the Krein form imposes neither. For a specified input preparation and phase-sensitive readout, the manuscript's Eq.~(95) bounds the transfer of the tested excitation by the coherence retained per prepared trial (Section~7.5). A measured coherence therefore limits the size of a violation for that preparation. It cannot establish that the transfer is exactly zero.

\paragraph{The ground.} $z_+$ is exact. Neither the vacuum alone nor the luminous ground alone is closed: $Q_A(f_0\oplus0)=z_+/\sqrt2$ and $Q_A(0\oplus g_1)=-z_+/\sqrt2$. The reduced space $\C g_0$ is orthogonal to both. Nor can the ground be kept in the reduced space by removing $z_+$ from the gauge. With $G'=G\ominus\C z_+$, the operator $P_{G'}J$ is again nilpotent and Krein self-adjoint, but its cohomology contains the hyperbolic plane spanned by $f_0$ and $g_1$. Its signature is $(1,2)$, a no-ghost violation that no sign convention removes. In the shift the ground event is exact, and removing it from the gauge makes the reduced space indefinite. More generally, a definite reduced space with the inherited form contains no nonzero null class, so a reduction that keeps the reduced space definite cannot retain the null ground event as a class. These statements about the shift use strong compatibility, $V^*y=x$. Under the weak bridge alone the relational ground of the manuscript need not be closed. On the unequal-cell path of the manuscript's Section~4, $\|Q_A(x\oplus y)/\sqrt2\|^2=\tfrac12(1-\sqrt{9/10})$, and the ground plane is not invariant under $Q_A$. The rank-one charge built from the restricted ground map $u_0$ of the manuscript's Section~4.4 is a valid operator in its own right, but it is not a restriction of $Q_A$ unless $V^*y=x$. What holds under R3 alone is that $y\in D$, so the reduced space $\ker A\oplus D^\perp$ is orthogonal to the luminous ground.

\section{What the reading claims, and what stays open}
\label{sec:reading}

\emph{Interpretation.} Two operations should be kept apart. If a physical state and its presentation are two descriptions of one event, identifying them leaves one copy of each event, with its matched probabilities. The BRST reduction does more: it removes the matched event altogether, and what survives is the appearance with no physical partner, evolving under its own Hamiltonian. The calculation is an exact account of what a specified presentation leaves unpaired. Under strong compatibility, $V^*y=x$, the anti-aligned companion of the manuscript's Section~14.2 is the non-closed partner of the aligned ground event. The premise that the observer's evolution acts on the quotient supplies the admissibility condition that the manuscript's conjecture lacked. A presentation is admissible when its polar map intertwines the dynamics, and a family is admissible when the polar map of its joint map does. These conditions are independent (Remark~\ref{rem:members}), and the exclusion of Corollary~\ref{cor:shift} holds under either. The premise concerns free evolution. It does not settle which interventions an observer can perform, and excluding transfers by interventions needs the second premise of Section~\ref{sec:shift}. The choice of the Krein form imposes neither premise. Whether an actual observer's evolution and interventions satisfy them is an empirical question. For a specified preparation and readout, Eq.~(95) bears on it.

\emph{The reading adopted here.} The relational event is \emph{neither one nor two}. It is one ray, the aligned graph vector, and its physical and phenomenal compressions are two distinct states that are not identified with each other. It is \emph{empty but not nonexistent}. As a vector of the carrier it has full Hilbert norm and carries the matched probabilities of the manuscript's Proposition~2.1. Its Krein norm is zero, and for a full-graph event its class in $H(Q_A)$ is zero. On this reading the Hilbert norm records that the seeing occurs, and the reduced space records what remains unpaired by the specified presentation. The reading is neither identification, which keeps one copy, nor erasure, which denies that the seeing occurs. A stronger reading would take the reduced space as the inventory of real appearances. We do not adopt it.

Two consequences follow. First, we propose to represent \emph{empty} by exactness, which is relative to a named charge, and not by Krein-nullness, which is a property of the vector alone. The two agree for full-graph events and come apart under the weak bridge. On the unequal-cell path the relational ground $(x\oplus y)/\sqrt2$ is Krein-null. It is exact for the rank-one charge of the restricted ground map $u_0=\ket y\!\bra x$, and for the full charge $Q_A$ it is not closed and has no class at all. A statement that something is empty must therefore name the charge. Second, a nonzero class marks a gap in what the specified presentation reaches. It confers no intrinsic existence on what fills the gap. At $\varepsilon=1$ the shift misses $g_0$: $A^*g_0=0$, so the direct return of $g_0$ through this presentation, $\|V^*g_0\|^2$, is zero. The shift rotated through $\pi/2$ sends $f_1$ to $-g_0$, and through it the direct return of $g_0$ is one. The same direction is unpaired relative to one presentation and paired relative to another. On the author's reading, this relational sense of emptiness is the Madhyamaka identification of emptiness with dependent arising. The calculation bears on Dharmak\=irti's criterion of reality, the capacity to produce an effect, only relative to a specified operation. The excess has zero direct return through this presentation. Its efficacy in another operation remains a separate question. The manuscript's other sense of emptiness, the vanishing of every grasper--grasped frame at the ground (its Sections~5.5 and 15.5), is a different property and should be kept distinct.

\emph{Commitments.}
\begin{itemize}
\item[(i)] Matched statistics are carried by the Hilbert norm and have no content in the reduced space. Paired effects are invariant but act as zero there (Corollary~\ref{cor:consequences}(a)), so the matched-probability correspondence of the manuscript's Proposition~2.1 becomes a statement about the gauge sector.
\item[(ii)] The reduced space is orthogonal to the luminous ground. Under strong compatibility the relational ground is exact, and keeping it in the reduced space makes that space indefinite: a definite reduced space contains no nonzero null class. Under the weak bridge alone it need not even be closed. A representation of the Third Yoga's ``luminosity still appears'' would have to live outside the reduced space, or the reading would have to change.
\item[(iii)] After all input redundancies of a family are treated as gauge, its reduced space equals the common mathematical defect $D_{\mathcal C}^\perp$ (sequel note, Corollary~3). For one preparation under R4 there are no redundancies, and $H(Q_A)=D^\perp$, which is zero when R5 fails. The conjecture of asymmetric nonduality claims more than a nonzero defect: a surviving direction must be a genuine possibility of appearance for an independently justified family. R4 on one preparation establishes neither the defect of a family nor that its directions are genuine possibilities of appearance. The manuscript's results that remove a common excess (a cyclic family of determinations, a second preparation setting, covariant channels) would leave the resolved reduced space zero.
\end{itemize}

\emph{Open.} Whether the Krein form should serve as an inner product for appearances is the departure from the manuscript and needs its own argument. Redundant presentations ($\ker B$) were handled here by restriction. The sequel note treats them systematically. Odd invariant operators, which exchange the sectors, were not treated. The shift is a toy model, and nothing here identifies an actual observer's preparations.

\end{document}
